
\begin{correctionbox}[Correction Exercice 1 : Loi Jointe et Marginales]
1.  On somme toutes les probabilités du tableau :
    $0.1 + 0.2 + 0.1 + 0.3 + 0.1 + 0.2 = 1.0$.
    Puisque la somme est 1 et toutes les probabilités sont non négatives, c'est une loi valide.

2.  Loi marginale de $X$ (somme des lignes) :
    $P(X=0) = P(X=0, Y=0) + P(X=0, Y=1) + P(X=0, Y=2) = 0.1 + 0.2 + 0.1 = 0.4$.
    $P(X=1) = P(X=1, Y=0) + P(X=1, Y=1) + P(X=1, Y=2) = 0.3 + 0.1 + 0.2 = 0.6$.

3.  Loi marginale de $Y$ (somme des colonnes) :
    $P(Y=0) = P(X=0, Y=0) + P(X=1, Y=0) = 0.1 + 0.3 = 0.4$.
    $P(Y=1) = P(X=0, Y=1) + P(X=1, Y=1) = 0.2 + 0.1 = 0.3$.
    $P(Y=2) = P(X=0, Y=2) + P(X=1, Y=2) = 0.1 + 0.2 = 0.3$.
\end{correctionbox}

\begin{correctionbox}[Correction Exercice 2 : Calcul de Probabilité Jointe]
1.  $P(X=0, Y \le 1) = P(X=0, Y=0) + P(X=0, Y=1) = 0.1 + 0.2 = 0.3$.
2.  $P(X=Y) = P(X=0, Y=0) + P(X=1, Y=1) = 0.1 + 0.1 = 0.2$.
3.  $P(X > Y) = P(X=1, Y=0) = 0.3$. (C'est la seule case où $x > y$).
\end{correctionbox}

\begin{correctionbox}[Correction Exercice 3 : Indépendance (Loi Jointe)]
On utilise les lois marginales de l'exercice 1 : $P(X=0)=0.4$ et $P(Y=0)=0.4$.
1.  $P(X=0) \times P(Y=0) = 0.4 \times 0.4 = 0.16$.
2.  Dans le tableau joint, $P(X=0, Y=0) = 0.1$.
3.  Puisque $P(X=0, Y=0) \neq P(X=0) \times P(Y=0)$ (car $0.1 \neq 0.16$), les variables $X$ et $Y$ \textbf{ne sont pas indépendantes}. (Un seul contre-exemple suffit).
\end{correctionbox}

% --- Corrections : Espérance, Covariance et Corrélation ---

\begin{correctionbox}[Correction Exercice 4 : Espérances Marginales]
1.  $E[X] = \sum_x x P(X=x) = (0)(P(X=0)) + (1)(P(X=1))$
    $E[X] = (0)(0.4) + (1)(0.6) = 0.6$.
2.  $E[Y] = \sum_y y P(Y=y) = (0)(P(Y=0)) + (1)(P(Y=1)) + (2)(P(Y=2))$
    $E[Y] = (0)(0.4) + (1)(0.3) + (2)(0.3) = 0 + 0.3 + 0.6 = 0.9$.
\end{correctionbox}

\begin{correctionbox}[Correction Exercice 5 : Espérance d'une Fonction (LOTUS)]
On somme $(xy)P(X=x, Y=y)$ sur les 6 cases. Les termes où $x=0$ ou $y=0$ sont nuls.
$E[XY] = (0 \cdot 0)(0.1) + (0 \cdot 1)(0.2) + (0 \cdot 2)(0.1) + (1 \cdot 0)(0.3) + (1 \cdot 1)(0.1) + (1 \cdot 2)(0.2)$
$E[XY] = 0 + 0 + 0 + 0 + (1)(0.1) + (2)(0.2) = 0.1 + 0.4 = 0.5$.
\end{correctionbox}

\begin{correctionbox}[Correction Exercice 6 : Covariance (Calcul)]
On utilise la formule $\text{Cov}(X,Y) = E[XY] - E[X]E[Y]$.
$$ \text{Cov}(X,Y) = 0.5 - (0.6)(0.9) = 0.5 - 0.54 = -0.04 $$
\end{correctionbox}

\begin{correctionbox}[Correction Exercice 7 : Variances Marginales]
1.  Pour $X$:
    $E[X^2] = (0^2)(0.4) + (1^2)(0.6) = 0.6$.
    $\text{Var}(X) = E[X^2] - (E[X])^2 = 0.6 - (0.6)^2 = 0.6 - 0.36 = 0.24$.
2.  Pour $Y$:
    $E[Y^2] = (0^2)(0.4) + (1^2)(0.3) + (2^2)(0.3) = 0 + 0.3 + (4)(0.3) = 0.3 + 1.2 = 1.5$.
    $\text{Var}(Y) = E[Y^2] - (E[Y])^2 = 1.5 - (0.9)^2 = 1.5 - 0.81 = 0.69$.
\end{correctionbox}

\begin{correctionbox}[Correction Exercice 8 : Corrélation (Calcul)]
On utilise la formule $\text{Corr}(X,Y) = \frac{\text{Cov}(X,Y)}{\sqrt{\text{Var}(X)\text{Var}(Y)}}$.
$$ \text{Corr}(X,Y) = \frac{-0.04}{\sqrt{0.24 \times 0.69}} = \frac{-0.04}{\sqrt{0.1656}} \approx \frac{-0.04}{0.4069} \approx -0.098 $$
La corrélation est très faible et négative.
\end{correctionbox}

% --- Corrections : Propriétés de la Variance et de la Covariance ---

\begin{correctionbox}[Correction Exercice 9 : Variance d'une Somme (Non Indépendant)]
On utilise la formule générale :
$$ \text{Var}(X+Y) = \text{Var}(X) + \text{Var}(Y) + 2\text{Cov}(X,Y) $$
$$ \text{Var}(X+Y) = 10 + 5 + 2(2) = 15 + 4 = 19 $$
\end{correctionbox}

\begin{correctionbox}[Correction Exercice 10 : Variance d'une Différence (Indépendant)]
1.  Puisque $X$ et $Y$ sont indépendantes, leur covariance est nulle : $\text{Cov}(X,Y) = 0$.
2.  On utilise la formule générale :
    $$ \text{Var}(X-Y) = \text{Var}(X + (-1)Y) = \text{Var}(X) + \text{Var}(-1 \cdot Y) + 2\text{Cov}(X, -Y) $$
    $$ = \text{Var}(X) + (-1)^2 \text{Var}(Y) - 2\text{Cov}(X, Y) $$
    $$ \text{Var}(X-Y) = \text{Var}(X) + \text{Var}(Y) - 2(0) $$
    $$ \text{Var}(X-Y) = 16 + 9 = 25 $$
\end{correctionbox}

\begin{correctionbox}[Correction Exercice 11 : Bilinéarité de la Covariance]
La covariance est linéaire sur son premier argument :
$$ \text{Cov}(X+Y, Z) = \text{Cov}(X, Z) + \text{Cov}(Y, Z) $$
\end{correctionbox}

\begin{correctionbox}[Correction Exercice 12 : Variance d'une Combinaison Linéaire]
On utilise $\text{Var}(aX + bY + c) = a^2 \text{Var}(X) + b^2 \text{Var}(Y) + 2ab\text{Cov}(X,Y)$.
Ici $a=3$, $b=-5$, $c=1$. $X$ et $Y$ sont indépendantes, donc $\text{Cov}(X,Y)=0$.
$$ \text{Var}(3X - 5Y + 1) = (3)^2 \text{Var}(X) + (-5)^2 \text{Var}(Y) + 0 $$
$$ = 9 \times (4) + 25 \times (2) = 36 + 50 = 86 $$
(Note : la constante additive $c=1$ ne change pas la variance).
\end{correctionbox}

\begin{correctionbox}[Correction Exercice 13 : Variance d'une Somme (Dés)]
1.  Oui, les lancers de deux dés standards sont des événements physiquement indépendants.
2.  Puisqu'ils sont indépendants, $\text{Cov}(D_1, D_2) = 0$.
    $$ \text{Var}(S) = \text{Var}(D_1 + D_2) = \text{Var}(D_1) + \text{Var}(D_2) $$
    $$ \text{Var}(S) = \frac{35}{12} + \frac{35}{12} = \frac{70}{12} = \frac{35}{6} $$
\end{correctionbox}

\begin{correctionbox}[Correction Exercice 14 : Covariance et Variance]
Par définition, $\text{Cov}(A, B) = E[(A-\mu_A)(B-\mu_B)]$.
Posons $A=X$ et $B=X$. Alors $\mu_A = \mu_X$ et $\mu_B = \mu_X$.
$$ \text{Cov}(X, X) = E[(X-\mu_X)(X-\mu_X)] = E[(X-\mu_X)^2] $$
C'est la définition de $\text{Var}(X)$.
\end{correctionbox}

\begin{correctionbox}[Correction Exercice 15 : Covariance avec une Constante]
On utilise la formule de calcul $\text{Cov}(X,c) = E[Xc] - E[X]E[c]$.
Par linéarité, $E[Xc] = cE[X]$.
L'espérance d'une constante est la constante elle-même : $E[c] = c$.
$$ \text{Cov}(X, c) = cE[X] - E[X]c = 0 $$
\end{correctionbox}

% --- Corrections : Standardisation et Somme de Poissons ---

\begin{correctionbox}[Correction Exercice 16 : Standardisation (Centrer-Réduire)]
$Z = \frac{X - 10}{2} = \frac{1}{2}X - 5$.
1.  Calcul de $E[Z]$ par linéarité :
    $$ E[Z] = E\left[ \frac{1}{2}X - 5 \right] = \frac{1}{2}E[X] - 5 = \frac{1}{2}(10) - 5 = 5 - 5 = 0 $$
2.  Calcul de $\text{Var}(Z)$ par les propriétés de la variance :
    $$ \text{Var}(Z) = \text{Var}\left( \frac{1}{2}X - 5 \right) = \left(\frac{1}{2}\right)^2 \text{Var}(X) = \frac{1}{4} \text{Var}(X) $$
    $$ \text{Var}(Z) = \frac{1}{4}(4) = 1 $$
    Par définition, une variable standardisée a une moyenne de 0 et une variance de 1.
\end{correctionbox}

\begin{correctionbox}[Correction Exercice 17 : Corrélation et Standardisation]
La corrélation $\text{Corr}(X,Y)$ EST, par définition, la covariance des versions standardisées $Z_X$ et $Z_Y$.
$$ \text{Cov}(Z_X, Z_Y) = \text{Corr}(X,Y) = 0.5 $$
\end{correctionbox}

\begin{correctionbox}[Correction Exercice 18 : Somme de Lois de Poisson]
1.  Puisque $X$ et $Y$ sont des v.a. de Poisson \textbf{indépendantes}, leur somme $S=X+Y$ suit aussi une \textbf{loi de Poisson}.
    Le nouveau paramètre est la somme des paramètres : $\lambda_S = \lambda_1 + \lambda_2 = 5 + 3 = 8$.
    Donc, $S \sim \text{Poisson}(\lambda=8)$.
2.  On cherche $P(S=6)$ pour $S \sim \text{Poisson}(8)$.
    $$ P(S=6) = \frac{e^{-8} 8^6}{6!} = \frac{e^{-8} \times 262144}{720} = 364.08 \times e^{-8} \approx 0.122 $$
\end{correctionbox}

\begin{correctionbox}[Correction Exercice 19 : Corrélation Nulle mais Dépendance]
1.  $E[X] = (-1)(1/3) + (0)(1/3) + (1)(1/3) = -1/3 + 0 + 1/3 = 0$.
2.  $E[XY] = E[X(X^2)] = E[X^3]$.
    $E[X^3] = (-1)^3(1/3) + (0)^3(1/3) + (1)^3(1/3) = (-1)(1/3) + 0 + (1)(1/3) = 0$.
3.  $\text{Cov}(X,Y) = E[XY] - E[X]E[Y] = 0 - (0)E[Y] = 0$.
    Les variables sont \textbf{non corrélées}.
4.  Les variables $X$ et $Y$ sont-elles indépendantes ? Non.
    Test : $P(X=1, Y=0) \stackrel{?}{=} P(X=1)P(Y=0)$.
    - $P(X=1, Y=0) = P(X=1, X^2=0) = 0$.
    - $P(X=1) = 1/3$.
    - $P(Y=0) = P(X^2=0) = P(X=0) = 1/3$.
    - $P(X=1)P(Y=0) = (1/3)(1/3) = 1/9$.
    Puisque $0 \neq 1/9$, elles \textbf{ne sont pas indépendantes}.
    C'est un exemple classique de dépendance non linéaire avec covariance nulle.
\end{correctionbox}

\begin{correctionbox}[Correction Exercice 20 : Bornes de la Corrélation]
$Y$ est une fonction linéaire parfaite de $X$ : $Y = aX + b$ avec $a=-3$ et $b=5$.
La corrélation $\text{Corr}(X,Y)$ mesure la force de la relation \textit{linéaire}. Puisqu'elle est parfaite, la corrélation doit être $\pm 1$.
Le coefficient $a = -3$ est négatif, donc la relation est décroissante.
Par conséquent, $\text{Corr}(X,Y) = -1$.
\end{correctionbox}
